%%%%%%%%%%%%%%%%%%%%%%%%%%%%%%%%%%%%%%%%%%%%%%%%%%%%%%%%%%%%%%%%%%%%%%
% Bizon -- DHBW deep learning server field manual
% Author: Prof. Dr.-Ing. Mark Schutera
% Based on Tufte-style layout
\makeindex
%%%%%%%%%%%%%%%%%%%%%%%%%%%%%%%%%%%%%%%%%%%%%%%%%%%%%%%%%%%%%%%%%%%%%%
%
% A field manual for the DHBW's Bizon deep learning server. Same shape as
% notes_dgxspark: short chapters, one topic each, every chapter also
% ships as its own handout PDF under chapter_pdfs/.
%
% Access follows the DHBW guide "Nutzung des Deep Learning Servers Bizon"
% (Anleitung fuer Studierende 1.2, bizon.bib): Sophos SSL VPN, then ssh
% with the DHBW account. Mark hands out only the setup package with the
% server address (student_setup/, gitignored, as for notes_dgxspark);
% neither the address nor any key material belongs in this book.
%
% Chapter spine.
%   01 Quickstart                  -- WRITTEN, from laptop to a shell on the Bizon
%   02 LLM gateway                 -- PLANNED, Open WebUI and LiteLLM API keys,
%                                     once GrazingBison runs on the Bizon
%
%%%%%%%%%%%%%%%%%%%%%%%%%%%%%%%%%%%%%%%%%%%%%%%%%%%%%%%%%%%%%%%%%%%%%%

\documentclass{../sharedAssets/tufte}

%%
% Book metadata
\title{Bizon}
\author[\defaultauthor]{\defaultauthor}
\publisher[\defaultpublisher]{\defaultpublisher}
\renewcommand{\notesdir}{notes_bizon}

% Load optional fonts
\loadoptionalfonts

% Additional packages
\usepackage{amsmath}
\usepackage{soul}
\usepackage{textcase}
\usepackage{enumitem}
\usepackage{tabularx}
\usepackage{fancyvrb}
\usepackage{float}
% Break long URLs anywhere so href'd links in the narrow Tufte margins
% wrap instead of running off the side of the page.
\usepackage{xurl}
\Urlmuskip=0mu plus 1mu\relax

% Let chapters open on any page instead of always on a recto. The shared
% class hardcodes \LoadClass[a5paper]{tufte-book} and forwards no options,
% so the flag is set here rather than passed (as in notes_3dprinting).
\makeatletter
\@openrightfalse
\makeatother

\begin{document}

\makestandardfrontmatter
\makededication{\openepigraph{``Freedom in a commons brings ruin to
all.''}{Garrett Hardin, The Tragedy of the Commons, 1968}}
\tableofcontents

% Introduction
\makeintroduction{%
  \newthought{This is a manual.} It gets you onto the DHBW's Bizon, a
  Linux deep learning server with NVIDIA GPUs, reachable from inside the
  DHBW network only.
  % TODO after bizon_recon.sh: GPU count, model, and memory per GPU here.
  This is a powerful compute substrate, use it. Then again you will never be
  alone on this machine, so be kind to data, jobs, and artifacts on this machine.
  Sandbox whatever you do, so you don't hamper the work of others.
  On the other hand, expect that someone else could nuke your stuff, so backups are advised.
  If you put stuff on that machine, expect that someone else will be able to see and use it.
  \vspace{1.5em}

  \marginnote{\newthought{The official documentation.} The DHBW guide
  for students, Nutzung des Deep Learning Servers Bizon, version 1.2. It
  also covers Miniconda, udocker containers, and Jupyter.}
  \newthought{This is currently a quickstart only,} so chapters on how to set up a proper project structure and environment
  are highly appreciated. Same holds for a chapter on how to effectively run jobs on a shared GPU server. Anything is worth adding that you would have loved to have known before diving deep.
  The goal is to have a short, concise, and practical manual for the Bizon over time.
  If a step is deprecated or you found a better way, contribute it back.
  }


%% Start the main matter
\mainmatter
\authorinrectoheader

% ----------------------------------------------------------------------
% Short chapters, one topic each.
% ----------------------------------------------------------------------
% !TEX root = ../notes_bizon.tex
% Chapter 1: Quickstart -- from your laptop to a shell on the Bizon
%%%%%%%%%%%%%%%%%%%%%%%%%%%%%%%%%%%%%%%%%%%%%%%%%%%%%%%%%%%%%%%%%%%%%%
%
% Access path: Sophos SSL VPN into the DHBW network, then ssh with the
% DHBW account, as in the DHBW guide (bizon.bib). Students log in as
% <username>@ab.ba.ba-ravensburg.de. The server address never appears in
% the book: it ships in the setup package (student_setup/, gitignored),
% handed out by mail as for the DGX Spark.
% TODO after bizon_recon.sh: confirm the ssh -l form and the GPU count.

\index{quickstart}\index{Bizon}
\chapter{Quickstart.}
\printchapterversion{01_quickstart}

\section{Where the buffalo roam}

\index{account}\newthought{The Bizon runs on your DHBW account.} Log in
with the username and password of the lab PCs. On the server, the
username carries the domain: \texttt{<username>@ab.ba.ba-ravensburg.de}
\cite{dhbw2021bizon}. The Bizon's address comes by mail, as a setup package.
\marginnote{\newthought{Scan the code} and send a mail stating your name,
course, and what you want to run.\par\medskip
\null\hfill\qrcode[height=1.0cm]{mailto:schutera@dhbw-ravensburg.de?subject=Bizon\%20Setup\&body=Name\%3A\%20\%0ACourse\%3A\%20\%0AProject\%3A\%20}}

\index{shared infrastructure}\newthought{Your account is one of many} on
a shared machine. Store nothing sensitive on it, leave other users' files
and processes alone, and sandbox your stuff.
Protect yourself from others, and protect others from yourself.
\marginnote{\newthought{Remember?} Chapter 1 of
Full Stack Handwerkszeug, Development and Production, covers this split,
along with SSH and Linux servers.
\url{https://material.schutera.com/lecturenotes/notes_missingsemester/notes_missingsemester.pdf}}%
The Bizon is used most efficiently when you spend as little time as possible to deploy your jobs and leave,
not for interactive development. If you need to develop interactively, please do so on your own machine.

\marginnote{\newthought{Scan the code} for the VPN user portal.\par\medskip
\null\hfill\qrcode[height=1.0cm]{https://vpn-rv.dhbw-ravensburg.de}}
\index{VPN}\newthought{The Sophos VPN takes you into the DHBW network.}
On campus you are already inside. Off campus, open
\url{https://vpn-rv.dhbw-ravensburg.de}, log in to the user portal with
your DHBW credentials, open the page \texttt{VPN}, and download the
client and the configuration. Then connect with your DHBW credentials:
\begin{itemize}[nosep]
  \item Windows: Sophos SSL VPN Client, tray icon, \texttt{Connect}.
  \item macOS: \texttt{.ovpn} into Tunnelblick \cite{tunnelblick2026}, \texttt{Connect}.
  \item Linux: \texttt{sudo openvpn -{}-config <file>.ovpn}.
\end{itemize}

\index{SSH}\newthought{Then log in} with the address from the package:
\begin{verbatim}
ssh -l <username>@ab.ba.ba-ravensburg.de <bizon-ip>
\end{verbatim}

\index{nvidia-smi}\newthought{Before you start a job,} run
\texttt{nvidia-smi} to see who is on which GPU. Pick a free one and pin
your job to it, for example GPU 3:
\begin{verbatim}
CUDA_VISIBLE_DEVICES=3 python train.py
\end{verbatim}


% Register the sources so the reference machinery has entries to resolve
% even while the book has a single chapter.
\nocite{dhbw2021bizon,tunnelblick2026}
\makestandardbackmatter{bizon}

\end{document}
